% !TeX program = xelatex
% arXiv TeX Live 2025 supports XeLaTeX, but not LuaLaTeX.
% xelatex -interaction=nonstopmode -halt-on-error -file-line-error noriega2026.tex
\documentclass[conference,twoside]{IEEEtran}

\usepackage{iftex}
\ifPDFTeX
  \PackageError{noriega2026}{Compile this manuscript with XeLaTeX}{The embedded PGFPlots typography requires fontspec/unicode-math so plot text uses Fira Sans and plot math uses Fira Math. pdfLaTeX cannot load OpenType Fira Math and will produce serif plot math.}
  \usepackage[T1]{fontenc}
  \usepackage{newpxtext}
  \usepackage{newpxmath}
  \usepackage{FiraSans}
  \usepackage[scaled=0.92]{beramono}
  \newcommand{\bodyfont}{\rmfamily}
  \newcommand{\figfont}{\sffamily}
  \DeclareRobustCommand{\plotmathversion}{}
\else
  \usepackage{fontspec}
  % Fonts are loaded by filename (resolved via kpathsea/TeX Live) rather than by
  % human name, so the document compiles on arXiv's XeLaTeX sandbox where the
  % system fontconfig database does not index TeX Live's bundled OpenType fonts.
  \setmainfont{texgyrepagella-regular.otf}[
    Ligatures=TeX,
    BoldFont=texgyrepagella-bold.otf,
    ItalicFont=texgyrepagella-italic.otf,
    BoldItalicFont=texgyrepagella-bolditalic.otf]
  \newfontfamily\bodyfont{texgyrepagella-regular.otf}[
    Ligatures=TeX,
    BoldFont=texgyrepagella-bold.otf,
    ItalicFont=texgyrepagella-italic.otf,
    BoldItalicFont=texgyrepagella-bolditalic.otf]
  \setsansfont{FiraSans-Regular.otf}[
    Ligatures=TeX,
    BoldFont=FiraSans-Bold.otf,
    ItalicFont=FiraSans-Italic.otf,
    BoldItalicFont=FiraSans-BoldItalic.otf]
  \newfontfamily\figfont{FiraSans-Regular.otf}[
    Ligatures=TeX,
    BoldFont=FiraSans-Bold.otf,
    ItalicFont=FiraSans-Italic.otf,
    BoldItalicFont=FiraSans-BoldItalic.otf]
  \setmonofont{FiraMono-Regular.otf}[
    BoldFont=FiraMono-Bold.otf,
    Scale=0.92]
  \usepackage{unicode-math}
  \setmathfont{texgyrepagella-math.otf}
  \setmathfont{FiraMath-Regular.otf}[version=plotsans]
  \DeclareRobustCommand{\plotmathversion}{\mathversion{plotsans}}
\fi

\let\Bbbk\relax
\usepackage{graphicx}
% svg package removed for arXiv submission (figures are pre-built PDFs)
\ifPDFTeX
  \usepackage{amsmath,amssymb}
\else
  \usepackage{amsmath}
\fi
\usepackage{booktabs}
\usepackage{tabularx}
\usepackage{array}
\usepackage{xcolor}
\usepackage{hyperref}
\hypersetup{
  colorlinks=true,
  linkcolor=blue!50!black,
  citecolor=blue!50!black,
  urlcolor=blue!50!black,
}
\urlstyle{tt}
\renewcommand{\UrlFont}{\ttfamily}
\DeclareRobustCommand{\monourl}[1]{{\ttfamily\href{#1}{\nolinkurl{#1}}}}
\usepackage{caption}
\usepackage{adjustbox}
\usepackage{placeins}
\usepackage{stfloats}
\usepackage{etoolbox}
\usepackage{tikz}
\usepackage{pgfplots}
\usepgfplotslibrary{groupplots,fillbetween}
\usetikzlibrary{calc}
\pgfplotsset{compat=1.18}
\usepackage{fancyhdr}

\definecolor{orcidgreen}{HTML}{A6CE39}
\DeclareRobustCommand{\orcidlink}[1]{%
  \href{https://orcid.org/#1}{%
    \tikz[baseline=-0.45ex]{%
      \fill[orcidgreen] (0,0) circle (0.82ex);
      \node[white,font=\sffamily\bfseries\tiny,inner sep=0pt] at (0,0) {iD};
    }%
  }%
}

% Thin rule above footnotes. This is deliberately shorter than float
% boundary rules so the two separators cannot be mistaken for each other.
\renewcommand{\footnoterule}{%
  \kern-3pt
  \hrule width 0.35\columnwidth height 0.4pt
  \kern 2.6pt
}

% Force every footnote to the same size and sans-serif face as captions.
\makeatletter
\long\def\@makefntext#1{%
  \parindent 1em\noindent
  {\normalfont\figfont\plotmathversion\scriptsize
   \hb@xt@1.8em{\hss\@makefnmark}#1}}
\makeatother

% Float separators are delegated to LaTeX's output routine. These hooks are
% called only after LaTeX has decided whether a float is top- or bottom-placed,
% which keeps the rules tied to the actual body/float transition. Top-placed
% tables carry their own booktabs bottom rule, so a table-only stack does not
% need another rule below it. Bottom-placed tables still need the boundary rule
% above the float stack, where body text transitions into the table.
\makeatletter
\newif\ifFloatStackHasFigure
\newif\ifFloatStackHasTable
\newcommand{\ScanFloatStackTypes}[1]{%
  \global\FloatStackHasFigurefalse
  \global\FloatStackHasTablefalse
  \begingroup
    \let\@elt\CheckFloatStackType
    #1%
  \endgroup
}
\newcommand{\CheckFloatStackType}[1]{%
  \@tempcnta\count#1\relax
  \divide\@tempcnta\@xxxii
  \ifnum\@tempcnta=\ftype@figure\relax
    \global\FloatStackHasFiguretrue
  \fi
  \ifnum\@tempcnta=\ftype@table\relax
    \global\FloatStackHasTabletrue
  \fi
}
\newcommand{\FloatBoundaryRule}{%
  \hrule width \linewidth height 0.4pt
}
\newcommand{\DoubleFloatTopBoundaryRule}{%
  \ifFloatStackHasFigure
    \vskip 2pt plus 0.5pt
    \hrule width \textwidth height 0.4pt
  \fi
}
\newcommand{\DoubleFloatBottomBoundaryRule}{%
  \ifFloatStackHasFigure
    \hrule width \textwidth height 0.4pt
    \vskip 2pt plus 0.5pt
  \else\ifFloatStackHasTable
    \hrule width \textwidth height 0.4pt
    \vskip 2pt plus 0.5pt
  \fi\fi
}
\newcommand{\BottomFloatBoundaryRule}{%
  \ifFloatStackHasFigure
    \FloatBoundaryRule
    \vskip 2pt plus 0.5pt
  \else\ifFloatStackHasTable
    \FloatBoundaryRule
    \vskip 2pt plus 0.5pt
  \fi\fi
}
\def\topfigrule{%
  \ifFloatStackHasFigure
    \vskip 2pt plus 0.5pt
    \FloatBoundaryRule
  \fi
}
\def\botfigrule{%
  \BottomFloatBoundaryRule
}
\def\dblfigrule{\hrule width \textwidth height 0.4pt}
\pretocmd{\@cflt}{\ScanFloatStackTypes\@toplist}{}{}
\pretocmd{\@cflb}{\ScanFloatStackTypes\@botlist}{}{}
\pretocmd{\@cdblflt}{\ScanFloatStackTypes\@dbltoplist}{}{}
\pretocmd{\@cdblflb}{\ScanFloatStackTypes\@dblbotlist}{}{}
\patchcmd{\@cdblflt}{\dblfigrule}{\DoubleFloatTopBoundaryRule}{}{}
\patchcmd{\@cdblflb}{\dblfigrule}{\DoubleFloatBottomBoundaryRule}{}{}
\makeatother
% Legacy calls are intentionally no-ops; separator placement is output-routine driven.
\newcommand{\floatcaprule}{}
\newcommand{\figcaprule}{}
\newcommand{\tabnoterule}{%
  \par\vspace{1pt}%
  \hrule width \linewidth height 0.4pt%
}
\renewcommand{\abstractname}{\bodyfont Abstract}
\makeatletter
\def\abstract{%
  \normalfont\bodyfont\bfseries\small
  \par\indent{\figfont\bfseries\itshape\abstractname}---\normalfont\bodyfont\bfseries\small\ignorespaces}
\def\endabstract{\par\normalfont\normalsize}
\makeatother

% Page-number headers
\pagestyle{fancy}
\fancyhf{}
% First page: centred page number at foot, no header
\fancypagestyle{plain}{%
  \fancyhf{}
  \fancyfoot[C]{\scriptsize\thepage}
  \renewcommand{\headrulewidth}{0pt}
}
% Subsequent pages: alternating author--page at outer margin
\fancyhead[RE]{\scriptsize\sffamily Noriega Cede\~{n}o,~p.\,\thepage}
\fancyhead[LO]{\scriptsize\sffamily Noriega Cede\~{n}o,~p.\,\thepage}
\renewcommand{\headrulewidth}{0pt}
\renewcommand{\footrulewidth}{0pt}

\pgfplotsset{
  colormap={viridisWarmSaturation}{
    rgb255(0cm)=(255,247,228);
    rgb255(25cm)=(248,213,167);
    rgb255(50cm)=(235,166,104);
    rgb255(75cm)=(216,112,55);
    rgb255(100cm)=(188,62,29)
  },
  colormap={orangeRedSequential}{
    rgb255(0cm)=(255,248,235);
    rgb255(25cm)=(252,210,151);
    rgb255(50cm)=(244,156,84);
    rgb255(75cm)=(216,89,54);
    rgb255(100cm)=(155,38,60)
  },
  colormap={viridisBlueGreen}{
    rgb255(0cm)=(247,252,250);
    rgb255(25cm)=(203,235,230);
    rgb255(50cm)=(132,205,198);
    rgb255(75cm)=(62,160,151);
    rgb255(100cm)=(33,145,140)
  },
}

\makeatletter
\renewcommand{\topfraction}{0.95}
\renewcommand{\bottomfraction}{0.90}
\renewcommand{\textfraction}{0.05}
\renewcommand{\floatpagefraction}{0.80}
\renewcommand{\dbltopfraction}{0.95}
\renewcommand{\dblfloatpagefraction}{0.80}
\setcounter{topnumber}{5}
\setcounter{bottomnumber}{5}
\setcounter{totalnumber}{10}
\setcounter{dbltopnumber}{4}
\setlength{\@fptop}{0pt}
\setlength{\@fpsep}{6pt plus 2pt minus 1pt}
\setlength{\@fpbot}{0pt plus 1pt}
\setlength{\@dblfptop}{0pt}
\setlength{\@dblfpsep}{6pt plus 2pt minus 1pt}
\setlength{\@dblfpbot}{0pt plus 1pt}
\makeatother
\setlength{\textfloatsep}{5pt plus 1pt minus 0.5pt}
\setlength{\floatsep}{6pt plus 2pt minus 1pt}
\setlength{\intextsep}{4pt plus 2pt minus 1pt}
\setlength{\dbltextfloatsep}{5pt plus 1pt minus 0.5pt}
\setlength{\dblfloatsep}{6pt plus 2pt minus 1pt}
\setlength{\abovecaptionskip}{2pt plus 0.5pt}
\setlength{\belowcaptionskip}{2pt plus 0.5pt}
\flushbottom

\DeclareCaptionFont{figsans}{\figfont\plotmathversion}
% Give figure bodies a little breathing room before the caption; keep tables tight.
\captionsetup{font={scriptsize,figsans},labelfont={bf,figsans}}
\captionsetup[figure]{skip=6pt}
\captionsetup[table]{skip=1pt}
\renewcommand{\arraystretch}{1.03}
\AtBeginDocument{%
  \setlength{\abovedisplayskip}{2pt plus 1pt minus 1pt}%
  \setlength{\belowdisplayskip}{2pt plus 1pt minus 1pt}%
  \setlength{\abovedisplayshortskip}{2pt plus 1pt minus 1pt}%
  \setlength{\belowdisplayshortskip}{2pt plus 1pt minus 1pt}%
  \setlength{\skip\footins}{8pt plus 2pt minus 1pt}%
  \setlength{\footnotesep}{3pt}%
}
\newcolumntype{Y}{>{\raggedright\arraybackslash}X}

\definecolor{condBaseline}{HTML}{444444}
\definecolor{condAugOne}{HTML}{7B1E3A}
\definecolor{condAugTwo}{HTML}{2D708E}
\definecolor{condFiltOne}{HTML}{20A387}
\definecolor{condFiltTwo}{HTML}{73D055}
\definecolor{classGlioma}{HTML}{440154}
\definecolor{classMeningioma}{HTML}{2D708E}
\definecolor{classNoTumour}{HTML}{20A387}
\definecolor{classPituitary}{HTML}{73D055}

\tikzset{
  boxAugOne/.style={draw=black, fill=condAugOne, line width=0.58pt},
  boxAugTwo/.style={draw=black, fill=condAugTwo, line width=0.58pt},
  boxFiltOne/.style={draw=black, fill=condFiltOne, line width=0.58pt},
  boxFiltTwo/.style={draw=black, fill=condFiltTwo, line width=0.58pt},
  boxSigAugOne/.style={draw=black, fill=condAugOne, line width=1.05pt},
  boxSigAugTwo/.style={draw=black, fill=condAugTwo, line width=1.05pt},
  boxSigFiltOne/.style={draw=black, fill=condFiltOne, line width=1.05pt},
  boxSigFiltTwo/.style={draw=black, fill=condFiltTwo, line width=1.05pt},
  whiskerAugOne/.style={draw=black, line width=0.48pt},
  whiskerAugTwo/.style={draw=black, line width=0.48pt},
  whiskerFiltOne/.style={draw=black, line width=0.48pt},
  whiskerFiltTwo/.style={draw=black, line width=0.48pt},
  whiskerSig/.style={draw=black, line width=0.58pt},
  meanAugOne/.style={draw=black, line width=0.72pt},
  meanAugTwo/.style={draw=black, line width=0.72pt},
  meanFiltOne/.style={draw=black, line width=0.72pt},
  meanFiltTwo/.style={draw=black, line width=0.72pt},
  meanSig/.style={draw=black, line width=0.72pt}
}

\newcommand{\plotfigurefont}{\figfont\plotmathversion}
\newcommand{\plotaxisfont}{\normalfont\plotfigurefont\scriptsize}
\newcommand{\plotaxisfontsmall}{\normalfont\plotfigurefont\scriptsize}
\newcommand{\plotaxistitlefont}{\normalfont\plotfigurefont\scriptsize}
\newcommand{\plotcellfont}{\normalfont\plotfigurefont\scriptsize}
\newcommand{\plotmatrixtitlefont}{\normalfont\plotfigurefont\scriptsize\bfseries}
\newcommand{\plottagfont}{\normalfont\plotfigurefont\scriptsize\bfseries}
\newcommand{\pgfplotlegendfont}{\normalfont\plotfigurefont\scriptsize}
\newcommand{\pct}[1]{#1\%}
\newcommand{\metriccell}[2]{#1~\textpm{}~#2}

\makeatletter
\def\section{\@startsection{section}{1}{\z@}{1.5ex plus 1.5ex minus 0.5ex}%
{0.7ex plus 1ex minus 0ex}{\normalfont\normalsize\figfont\scshape\bfseries\centering}}
\def\subsection{\@startsection{subsection}{2}{\z@}{1.5ex plus 1.5ex minus 0.5ex}%
{0.7ex plus .5ex minus 0ex}{\normalfont\normalsize\bodyfont\bfseries\raggedright}}
\def\subsubsection{\@startsection{subsubsection}{3}{\z@}{1.5ex plus 1.5ex minus 0.5ex}%
{0.5ex plus .5ex minus 0ex}{\normalfont\normalsize\bodyfont\raggedright}}
\def\paragraph{\@startsection{paragraph}{4}{2\parindent}{0ex plus 0.1ex minus 0.1ex}%
{0ex}{\normalfont\normalsize\bodyfont\itshape}}
\def\subparagraph{\@startsection{subparagraph}{5}{2\parindent}{0ex plus 0.1ex minus 0.1ex}%
{0ex}{\normalfont\normalsize\bodyfont\itshape}}
\makeatother

\AtBeginDocument{\bodyfont}
% Standalone PGFPlots source for the corresponding pre-built figure PDF.
% The main arXiv manuscript includes figures/pgfplots_prebuilt/*.pdf to avoid
% TeX memory exhaustion while preserving the exact embedded-plot typography.
\begin{document}
\thispagestyle{empty}
\noindent
% --- Embedded source from figures/pgfplots/pgfplots_p020_p021/P020_P021_cnn_v2_confusion_matrices.tex ---
\begingroup\plotfigurefont
\newcommand{\figtext}{\plotcellfont}

\pgfplotsset{
  colormap={wineRedSequential}{
    rgb255(0cm)=(255,247,249);
    rgb255(25cm)=(239,203,213);
    rgb255(50cm)=(210,139,158);
    rgb255(75cm)=(164,76,103);
    rgb255(100cm)=(123,30,58)
  },
  colormap={viridisBlueGreen}{
    rgb255(0cm)=(247,252,250);
    rgb255(25cm)=(203,235,230);
    rgb255(50cm)=(132,205,198);
    rgb255(75cm)=(62,160,151);
    rgb255(100cm)=(33,145,140)
  },
  every axis/.append style={
    scale only axis,
    axis on top=true,
    axis line style={black, line width=0.55pt},
    tick style={black, line width=0.55pt},
    tick align=outside,
    tick pos=left,
    major tick length=2.2pt,
    title style={font=\plotmatrixtitlefont},
    label style={font=\plotaxistitlefont},
    tick label style={font=\plotaxisfont, /pgf/number format/assume math mode=true},
  },
}

\begin{adjustbox}{max width=\linewidth,center}
\begin{tikzpicture}
\begin{axis}[
  name=tumour,
  title={(a)},
  width=53mm,
  height=53mm,
  xmin=0.5,
  xmax=4.5,
  ymin=0.5,
  ymax=4.5,
  axis equal image,
  enlargelimits=false,
  xtick={1,2,3,4},
  ytick={4,3,2,1},
  xticklabels={Glioma,Mening.,No tumour,Pituitary},
  yticklabels={Glioma,Mening.,No tumour,Pituitary},
  x tick label style={font=\plotaxisfont, /pgf/number format/assume math mode=true, rotate=35, anchor=east},
  y tick label style={font=\plotaxisfont, /pgf/number format/assume math mode=true, rotate=35, anchor=east},
  xlabel={},
  ylabel={},
  point meta min=0,
  point meta max=100,
  colormap name=wineRedSequential,
  colorbar,
  colorbar style={
    height=53mm,
    ylabel={},
    ytick={0,25,50,75,100},
    ytick pos=right,
    axis line style={black, line width=0.55pt},
    tick style={black, line width=0.55pt},
    tick align=outside,
    major tick length=2.2pt,
    scaled y ticks=false,
    yticklabel style={font=\plotaxisfont},
    ylabel style={font=\plotaxistitlefont}
  },
]
\addplot[
  matrix plot*,
  mesh/cols=4,
  point meta=explicit,
  draw=white,
  line width=0.3pt
] table[x=x, y=y, meta=mean] {../pgfplots/pgfplots_p020_p021/data/tumour_mean_confusion.dat};
\draw[black, line width=0.55pt] (axis cs:0.5,0.5) rectangle (axis cs:4.5,4.5);
\node[align=center, font=\figtext, text=white, inner sep=0pt] at (axis cs:1,4) {96.2 $\pm$ 0.7};
\node[align=center, font=\figtext, text=black, inner sep=0pt] at (axis cs:2,4) {3.6 $\pm$ 0.7};
\node[align=center, font=\figtext, text=black, inner sep=0pt] at (axis cs:3,4) {0.1 $\pm$ 0.0};
\node[align=center, font=\figtext, text=black, inner sep=0pt] at (axis cs:4,4) {0.0 $\pm$ 0.0};
\node[align=center, font=\figtext, text=black, inner sep=0pt] at (axis cs:1,3) {3.9 $\pm$ 0.9};
\node[align=center, font=\figtext, text=white, inner sep=0pt] at (axis cs:2,3) {87.4 $\pm$ 1.8};
\node[align=center, font=\figtext, text=black, inner sep=0pt] at (axis cs:3,3) {6.4 $\pm$ 2.4};
\node[align=center, font=\figtext, text=black, inner sep=0pt] at (axis cs:4,3) {2.3 $\pm$ 0.2};
\node[align=center, font=\figtext, text=black, inner sep=0pt] at (axis cs:1,2) {0.0 $\pm$ 0.0};
\node[align=center, font=\figtext, text=black, inner sep=0pt] at (axis cs:2,2) {0.0 $\pm$ 0.0};
\node[align=center, font=\figtext, text=white, inner sep=0pt] at (axis cs:3,2) {100.0 $\pm$ 0.0};
\node[align=center, font=\figtext, text=black, inner sep=0pt] at (axis cs:4,2) {0.0 $\pm$ 0.0};
\node[align=center, font=\figtext, text=black, inner sep=0pt] at (axis cs:1,1) {0.1 $\pm$ 0.1};
\node[align=center, font=\figtext, text=black, inner sep=0pt] at (axis cs:2,1) {0.9 $\pm$ 0.2};
\node[align=center, font=\figtext, text=black, inner sep=0pt] at (axis cs:3,1) {0.0 $\pm$ 0.0};
\node[align=center, font=\figtext, text=white, inner sep=0pt] at (axis cs:4,1) {98.9 $\pm$ 0.3};
\end{axis}

\begin{axis}[
  name=plane,
  at={(tumour.south west)},
  anchor=north west,
  yshift=-24mm,
  title={(b)},
  width=53mm,
  height=53mm,
  xmin=0.5,
  xmax=3.5,
  ymin=0.5,
  ymax=3.5,
  axis equal image,
  enlargelimits=false,
  xtick={1,2,3},
  ytick={3,2,1},
  xticklabels={Axial,Coronal,Sagittal},
  yticklabels={Axial,Coronal,Sagittal},
  x tick label style={font=\plotaxisfont, /pgf/number format/assume math mode=true, rotate=35, anchor=east},
  y tick label style={font=\plotaxisfont, /pgf/number format/assume math mode=true, rotate=35, anchor=east},
  xlabel={},
  ylabel={},
  point meta min=0,
  point meta max=100,
  colormap name=viridisBlueGreen,
  colorbar,
  colorbar style={
    height=53mm,
    ylabel={},
    ytick={0,25,50,75,100},
    ytick pos=right,
    axis line style={black, line width=0.55pt},
    tick style={black, line width=0.55pt},
    tick align=outside,
    major tick length=2.2pt,
    scaled y ticks=false,
    yticklabel style={font=\plotaxisfont},
    ylabel style={font=\plotaxistitlefont}
  },
]
\addplot[
  matrix plot*,
  mesh/cols=3,
  point meta=explicit,
  draw=white,
  line width=0.3pt
] table[x=x, y=y, meta=mean] {../pgfplots/pgfplots_p020_p021/data/plane_mean_confusion.dat};
\draw[black, line width=0.55pt] (axis cs:0.5,0.5) rectangle (axis cs:3.5,3.5);
\node[align=center, font=\figtext, text=white, inner sep=0pt] at (axis cs:1,3) {99.3 $\pm$ 0.2};
\node[align=center, font=\figtext, text=black, inner sep=0pt] at (axis cs:2,3) {0.6 $\pm$ 0.1};
\node[align=center, font=\figtext, text=black, inner sep=0pt] at (axis cs:3,3) {0.1 $\pm$ 0.1};
\node[align=center, font=\figtext, text=black, inner sep=0pt] at (axis cs:1,2) {1.7 $\pm$ 0.2};
\node[align=center, font=\figtext, text=white, inner sep=0pt] at (axis cs:2,2) {96.5 $\pm$ 0.6};
\node[align=center, font=\figtext, text=black, inner sep=0pt] at (axis cs:3,2) {1.9 $\pm$ 0.5};
\node[align=center, font=\figtext, text=black, inner sep=0pt] at (axis cs:1,1) {0.1 $\pm$ 0.0};
\node[align=center, font=\figtext, text=black, inner sep=0pt] at (axis cs:2,1) {1.5 $\pm$ 0.1};
\node[align=center, font=\figtext, text=white, inner sep=0pt] at (axis cs:3,1) {98.4 $\pm$ 0.1};
\end{axis}
\node[font=\figtext, rotate=90, anchor=center] at
  ($ (tumour.outer north west)!0.5!(plane.outer south west) + (4mm,0) $)
  {True};
\node[font=\figtext, rotate=90, anchor=center] at
  ($ (tumour.outer north east)!0.5!(plane.outer south east) + (21mm,0) $)
  {Mean row-normalized (\%)};
\node[font=\figtext, anchor=north] at
  ($ (plane.south west)!0.5!(plane.south east) + (0,-12mm) $)
  {Predicted};
\end{tikzpicture}%
\end{adjustbox}
\endgroup

\end{document}
